% ============================================================================
\section{Distillation and Task Allocation}
\label{sec:setting}
% ============================================================================

\subsection{Dense supervision on student rollouts}

There are $M$ tasks with prompt distributions $\mathcal D_i$ and assigned
frozen teachers $\pi_{T_i}$. The student and teachers share a vocabulary.
For a task-$i$ prompt $x$, the student generates a response $y$. At position
$t$, the teacher scores the student-generated prefix $c=(x,y_{<t})$.
Write $p_\theta(a\mid c)=\pi_\theta(a\mid c)$ and
$q_i(a\mid c)=\pi_{T_i}(a\mid c)$.

We use the teacher top-$k$ loss from MOPD
\citep{ma2026mopdmultiteacheronpolicydistillation}. Let
$S_i(c)=\operatorname{TopK}_k(q_i(\cdot\mid c))$. The position loss is
\begin{equation}
    \ell_i^{(k)}(\theta;c)=
    \sum_{a\in S_i(c)}
    \left[
        p_\theta(a\mid c)\log\frac{p_\theta(a\mid c)}{q_i(a\mid c)}
        -p_\theta(a\mid c)+q_i(a\mid c)
    \right].
    \label{eq:dense_topk_loss}
\end{equation}
The probabilities retain their full-vocabulary normalization. The
$-p_\theta+q_i$ correction makes the loss smallest when student and teacher
probabilities agree on the retained tokens. Both the teacher scores and the
generated prefixes are held fixed during the update. This retained-support objective gives a dense gradient at each
observed position using a small teacher payload.
GPAS operates on the resulting micro-batch gradients; the allocation
method also applies to other distillation losses.

We average position losses within each response, then average responses
within a task. Positive weights $w_i$, fixed before training and summing
to one, specify the desired balance:
\begin{equation}
    F(\theta)=\sum_{i=1}^M w_iL_i(\theta).
    \label{eq:fixed_multitask_objective}
\end{equation}
Here $L_i$ is the expected response-mean loss under the current rollout
distribution. Each update differentiates the loss on those sampled prefixes;
Appendix~\ref{app:local_objective} gives the corresponding step-local
objective. Our main study uses equal task weights.

\subsection{Averaging task gradients before combining them}

A micro-batch contains $b$ prompts from one task. An optimizer step processes
$G$ micro-batches, with task $i$ receiving $m_i$ of them. The counts sum to
$G$ and satisfy $m_{\min}\leq m_i\leq m_{\max}$. We use
$m_{\min}\geq2$ to estimate each task's variability from the same step.

Let $g_{i,s}$ be the gradient of the response-mean loss in micro-batch $s$
of task $i$. GPAS accumulates
\begin{equation}
    A=\sum_{i=1}^M w_i
       \underbrace{\frac1{m_i}\sum_{s=1}^{m_i}g_{i,s}}_{\text{task mean}}.
    \label{eq:batch_gradient_observation}
\end{equation}
Implementation only requires multiplying each micro-batch loss by $w_i/m_i$.
For example, an equally weighted task contributes one quarter of the
combined mean whether it receives two or six micro-batches; six
micro-batches give a more precise estimate of that mean.

Counts are chosen before drawing the step's data. If fresh task-$i$
micro-batches have mean $\mu_i$ at the current training state, then
$\mathbb E[A\mid m]=\sum_iw_i\mu_i$. Thus, the scheduler can change counts
without changing the intended weighted mean. Response-level averaging also
prevents a task from gaining extra loss weight simply because its responses
are longer. This addresses the length imbalance highlighted by Open-MOPD
\citep{gao2026openmopddiagnosingfixingcapability}, while making the response
the unit of averaging within each task.
